\chapter[Diseña adecuadamente según estándares de calidad][Diseña según estándares de calidad]{Diseña adecuadamente según estándares de calidad}
\label{cap:diseno}

Este capítulo presenta el diseño del instrumento: estándares, seis figuras de cuerpo y un cuadro de mecanismos, las decisiones
que las explican con su razón y su evidencia, los atributos de calidad que persigue y el diseño
que las figuras no sustituyen. No presenta resultados: las cifras que lo evalúan pertenecen a
las etapas 3 y 4.

\section{Estándares y enfoque}
\label{sec:estandares}

El diseño se describe desde varios puntos de vista, cada uno dirigido a una preocupación de
las partes interesadas, como prescriben ISO/IEC/IEEE 42010 \parencite{iso42010}, IEEE 1016
\parencite{ieee1016} y SWEBOK \parencite[cap.~2, §2.1]{swebok2025}, con la separación en vistas
de \textcite{kruchten1995views}. Cada vista responde una pregunta, enunciada en el párrafo que
la introduce; su pie da el mensaje y la clave al pie de la figura define toda marca que no sea
UML. Las vistas usan UML 2.5.1 \parencite{omg_uml251} como notación estructural y de
comportamiento \parencite[cap.~11, §2]{swebok2025}: contexto (F1), casos de uso (F2), modelo conceptual
(F6), clases (A1), actividad (F4, F7) y secuencia (A2). F3 es una correspondencia de
responsabilidades y T5 un cuadro de mecanismos; ninguno pretende expresar dependencias UML. Cinco
principios gobiernan el sistema: una tubería por lotes sobre archivos \parencite[cap.~2,
§2.2]{swebok2025}; la separación de responsabilidades orientada al ocultamiento de información \parencite{parnas1972criteria} en el motor
fiel \parencite[cap.~3, §1.4]{swebok2025}; la concurrencia y el manejo de errores como
decisiones explícitas \parencite[cap.~3, §3.1 y §3.5]{swebok2025}; el motor moderno como
variante de una línea de productos \parencite[cap.~3, §3.8]{swebok2025}; y la verificación
como parte del sistema, con la referencia como oráculo \parencite[cap.~5, §1.2.7]{swebok2025}.
Cada decisión que explica una vista se cita como D-xx, con la numeración del registro de la
tesis (D-01 a D-21) continuada desde D-22, y se reúne con su razón y su evidencia en la sección
\ref{sec:decisiones}, como pide SWEBOK \parencite[cap.~3, §4.6]{swebok2025}. Cada elemento
dibujado tiene una cita de archivo y línea (\href{run:../figuras/FACTS.md}{juego de auditoría});
la notación está en el anexo \ref{anexo:figuras}.

\section{Vistas del instrumento y mecanismos numéricos}
\label{sec:vistas}

\subsection{Contexto y alcance del instrumento}
\label{sec:F1}

¿Quién usa el instrumento, con qué entradas y para obtener qué? En la figura \ref{fig:F1}, el
analista declara un panel, semillas y el motor, con lo que ejecuta una variante identificada
(cuadro \ref{tab:referentes}), y recibe coordenadas y parámetros; el verificador solicita
comparaciones. Son externos las fuentes de paneles (VoteView \parencite{boche2018voteview}; la
base de QueVotan y los datos de \textcite{fabrega2025ideological}), el servicio de semillas
\parencite{wnominate_pkg} y la referencia wmay, que devuelve coordenadas o una
log-verosimilitud. El portal en R es la entrada común a los dos motores y a la referencia,
cuyos binarios también pueden ejecutarse directamente; aplica el linaje del motor fiel y sella
cada corrida aceptada (D-24, D-30). La tubería es por lotes sobre archivos porque la referencia solo lee y escribe
archivos; la integración con QueVotan (anexo \ref{anexo:contratos}) es externa a este
trabajo.

\begin{figure}[tbp]
\centering
\makebox[\textwidth][c]{\includegraphics[width=13.5cm]{F1-contexto.pdf}}
\caption[F1. Contexto del instrumento]{Contexto lógico del instrumento de investigación. Ruta común: portal R, ejecutables y validación de archivos (D-24, D-30).}
\label{fig:F1}
\end{figure}

\subsection{Objetivos de los usuarios}
\label{sec:F2}

¿Qué puede hacer cada rol con el instrumento? En la figura \ref{fig:F2}, el analista declara
panel y semillas y estima posiciones, caso abstracto que se especializa en estático (RF03) y
dinámico (RF04), con el motor fiel o el moderno (RF05, RF06); la carga y validación de
entradas (RF01) ocurre dentro de la estimación. Validar con el caso chileno (RF10 a RF12) es
un flujo de investigación que reutiliza la estimación y no una función de la línea de
comandos. El sujeto incluye estimación, validación y comparación, y la elipse está dentro;
la reformulación del objetivo chileno sigue siendo propuesta de la sección
\ref{sec:propuesta-chile}. El verificador compara contra wmay (RF08, RF09); ambos roles pueden
ser la misma persona. Las narrativas de cada caso están en el anexo \ref{anexo:narrativas}.

\begin{figure}[tbp]
\centering
\makebox[\textwidth][c]{\includegraphics[width=13.5cm]{F2-casos-de-uso.pdf}}
\caption[F2. Objetivos de los usuarios]{Casos de uso del instrumento de investigación, que comprende estimación y validación. RF remite a los requisitos vigentes; la propuesta sobre Chile sigue pendiente.}
\label{fig:F2}
\end{figure}

\subsection{Del Fortran al C++: organización}
\label{sec:F3}

¿Dónde quedó cada responsabilidad del instrumento heredado? La figura \ref{fig:F3}
corresponde responsabilidades, sin superponer el grafo de dependencias. Distingue el programa
con archivos y estado \texttt{COMMON} de 2004 de la subrutina \texttt{dwnom} con módulos y
envoltorios de wmay. La derecha identifica entrada, configuración explícita, orquestación,
búsquedas, geometría, verosimilitud y soporte en el motor fiel. El catálogo de rutinas y el
grafo de inclusiones se documentan en el anexo \ref{anexo:correspondencia}.

El estado sigue concentrado en DWNominate para permitir el volcado por fases (D-22); las
búsquedas son funciones libres (D-23). \texttt{GRID} queda tras un interruptor (D-24) y el
perfil comparado usa \texttt{DGESDD} (D-01). Esa organización persigue mantenibilidad
\parencite{parnas1972criteria}; no demuestra ocultamiento completo ni una mejora medida.

\begin{figure}[tbp]
\centering
\makebox[\textwidth][c]{\includegraphics[width=13.5cm]{F3-fortran-a-cpp.pdf}}
\caption[F3. Correspondencia de responsabilidades]{Correspondencia de responsabilidades entre original 2004, wmay y fiel C++; catálogo de rutinas en el anexo \ref{anexo:correspondencia}.}
\label{fig:F3}
\end{figure}

\subsection{Ciclo de estimación compartido}
\label{sec:F4}

La figura \ref{fig:F4} ordena las actualizaciones de parámetros, votaciones y
legisladores, con evaluaciones intermedias. Tras cargar panel y semillas, la validez exige
un margen de minoría de al menos $2{,}5$\,\%. El corte usa \texttt{CUTPLANE} en 2D o
\texttt{JAN11PT} en 1D, seguido de \texttt{RCINT2}. El perfil wmay/fiel repite $T$ ciclos
fijos y exporta parámetros terminales (D-07); el resumen fiel mezcla fases (M-02, I2).
No hay prueba de convergencia. Compartir este flujo
no garantiza cifras iguales: la verosimilitud no es cóncava y se compara a ciclos igualados
(algoritmo \ref{alg:ciclo}; anexo \ref{anexo:perfil}).

\paragraph{Concurrencia.} El motor fiel usa OpenMP en votaciones, legisladores y la
verosimilitud de las búsquedas globales; la evaluación entre fases es serial; la referencia ejecuta todas las fases secuencialmente
\parencite[cap.~3, §3.1]{swebok2025}. Para conciliar paralelismo y determinismo (RNF02),
ninguna cantidad de punto flotante se reduce en una región paralela: la verosimilitud se
acumula por legislador y se suma después en orden de índice (D-03). Los hilos los fija
\archivo{OMP_NUM_THREADS}, que el portal impone, y toda comparación usa uno (D-28). La fase de
legisladores inserta en un contenedor compartido sin sección crítica (P-01); se corrige con la
preinicialización en serie que ya protege a los coeficientes temporales y, hasta entonces, no
se afirma ninguna propiedad del motor fiel con más de un hilo.

\begin{figure}[tbp]
\centering
\makebox[\textwidth][c]{\includegraphics[width=13.5cm]{F4-ciclo-estimacion.pdf}}
\caption[F4. Actividad de estimación]{Actividad del perfil wmay/fiel: 2D, cuatro ciclos, parámetros terminales y resumen fiel mixto (M-02). No establece equivalencia con el original 2004.}
\label{fig:F4}
\end{figure}

\subsection{Intervención moderna}
\label{sec:T5}

¿Qué búsqueda cambia en cada actualización? El cuadro \ref{tab:T5} compara mecanismos
de implementación. El moderno conserva el modelo temporal y reemplaza búsquedas heredadas
por BOBYQA escalar \parencite{powell2009bobyqa} y COBYLA de bloques
\parencite{powell1994cobyla}, mediante NLopt \parencite{johnson_nlopt}. El perfil seleccionado
usa intervalos escalares locales y restricciones de norma; SLSQP y estrategias híbridas son
opciones fuera de ese perfil. Son implementaciones separadas (D-11, D-12).

El perfil de replicación fija la tabla de verosimilitud de la referencia y retira las cotas
globales de $w_2$ y $\beta$; persiste el intervalo local, y el portal añade precisión estricta
(D-13, D-25). La aceptación de bloques y el control de descenso global también son mecanismos
propios del moderno (D-32): el perfil no vuelve idénticos los recorridos de búsqueda. Panel,
semillas, ciclos y evaluación son controles del experimento, descritos en el cuadro
\ref{tab:protocolos} y el anexo \ref{anexo:perfil}. Esta intervención numérica es distinta
de la futura optimización de memoria con salidas idénticas; sus diferencias no se califican
automáticamente como errores de traducción.

\begin{table}[tbp]
\centering
\makebox[\textwidth][c]{\includegraphics[width=13.5cm]{T5-mecanismos-modernos.pdf}}
\caption[T5. Mecanismos de la intervención moderna]{Mecanismos fieles y modernos en el perfil
controlado, verificados contra la implementación (D-13, D-25, D-32).}
\label{tab:T5}
\end{table}

\subsection{Datos y tiempo servido}
\label{sec:F6}

¿Qué significa cada fila y cómo se reconstruyen sus coordenadas? En la figura \ref{fig:F6}, un
legislador tiene una fila servida por periodo presente, con sus celdas de voto. Hay dos
contratos de semilla: el archivo por legislador, que el binario exige aunque cada legislador
pueda faltar en él, y el archivo por fila, opcional y con prioridad; sin ninguna, el
legislador parte del origen, como en la referencia, porque el origen es neutral (D-26). El
portal rechaza sembrar en proceso un panel de varios periodos, que dejaría en el origen a
quien no sirvió en el primero (D-27), y la comparación con wmay usa la semilla por
legislador, único contrato que esa referencia acepta. La trayectoria se evalúa solo en los
periodos servidos, con un tiempo local, y su grado efectivo es el de la regla de la referencia
(D-06): por eso el legislador del ejemplo, con dos periodos, recibe las mismas coordenadas en
sus dos filas. Como la escala de tiempo es propia de cada legislador, se exportan las
coordenadas en los periodos servidos y no sobre una rejilla global (D-10).

\begin{figure}[tbp]
\centering
\makebox[\textwidth][c]{\includegraphics[width=13.5cm]{F6-datos-tiempo-servido.pdf}}
\caption[F6. Datos y tiempo servido]{Modelo conceptual de panel, filas servidas, semillas y trayectoria. El grado efectivo depende del número de periodos servidos (D-06, D-10, D-27).}
\label{fig:F6}
\end{figure}

\subsection{Proceso de verificación}
\label{sec:F7}

¿Qué comprobación permite confiar en la reimplementación y qué limita esa conclusión? La
figura \ref{fig:F7} muestra dos de los tres niveles de RF08 sobre corridas controladas, con el
mismo panel, la semilla por legislador y el mismo horizonte; el de componente está en el
cuadro \ref{tab:niveles}. (a) El estado terminal C++ se evalúa con la verosimilitud de wmay,
con cero ciclos, y se compara con la que la corrida wmay escribe en su resumen. (b) Las
coordenadas se emparejan por (legislador, periodo), se centran y se alinean por una
transformación ortogonal, rotación o reflexión, sin reescalar, y se informan las filas
emparejadas, las correlaciones por dimensión $r_1$ y $r_2$, la razón de amplitud $A$ (ecuación
\ref{eq:amplitud}) y la distancia media, con cobertura y marco. Es verificación, no validación
\parencite{oberkampf2010verification}; la amplitud acompaña a la correlación, invariante a la
escala \parencite{herron2004scale}, y el informe (RF09) añade el desacuerdo por votación
(D-14). Ninguna operación demuestra igualdad ni reemplaza
el control de paridad de entradas, documentado en el anexo \ref{anexo:perfil}; esta vista
usa wmay. El protocolo de aceptación de la tesis se declara por separado a continuación.

\begin{figure}[tbp]
\centering
\makebox[\textwidth][c]{\includegraphics[width=13.5cm]{F7-verificacion.pdf}}
\caption[F7. Verificación con wmay]{Puntuación de estado y comparación de mapas del protocolo wmay. Cobertura y marco explícitos, alineación ortogonal sin reescalar (D-14).}
\label{fig:F7}
\end{figure}

\section{Decisiones de diseño y su razón}
\label{sec:decisiones}

El cuadro \ref{tab:decisiones} reúne las decisiones que explican las vistas, con la
alternativa descartada, la razón y la línea de código o configuración que la sostiene,
releída para este informe; las existentes conservan su número de la tesis y las nuevas siguen
desde D-22. El registro completo está en el anexo \ref{anexo:decisiones}, y la sección
\ref{sec:registros} lo agrupa en registros de decisión. Las rutas remiten al repositorio del
motor (\href{run:../evidence/engine/README.md}{motores, revisión ad945a9}) y al portal
(\href{run:../evidence/research/reproduce/R/dwnom.R}{portal en R}).

{\scriptsize
\setlength{\tabcolsep}{2pt}
\begin{longtable}{L{0.85cm}L{3.1cm}L{2.2cm}L{3.2cm}L{2.25cm}L{0.95cm}}
\caption[Decisiones de diseño que explican las vistas.]{Decisiones de diseño que explican las
vistas. Evidencia: F, \archivo{engines/faithful}; M, \archivo{engines/modern}; R,
\archivo{reproduce/R/dwnom.R}, con números de línea.}\label{tab:decisiones}\\
\toprule
Id & Decisión & Alternativa descartada & Razón & Evidencia & Figura \\
\midrule
\endfirsthead
\multicolumn{6}{l}{\scriptsize\emph{Cuadro \ref{tab:decisiones}, continuación}}\\
\toprule
Id & Decisión & Alternativa descartada & Razón & Evidencia & Figura \\
\midrule
\endhead
\bottomrule
\endfoot
D-22 & estado en \texttt{DWNominate}, trazado miembro a miembro a \texttt{xxcom\_mod}; fases
como métodos en orden & objetos por fase & la verificación por componente detiene el motor
entre fases y compara arreglos con la referencia & F \archivo{dwnominate.hpp:5-8} & F3, A1 \\
D-23 & búsquedas como funciones libres, contexto por argumento & métodos del orquestador & una
búsqueda sin estado propio se prueba aislada & F \archivo{grid_optimizer.hpp:126-151}; M
\archivo{tests/} & F3, A1 \\
D-24 & \texttt{GRID} de 2004 tras un interruptor que apaga \archivo{--wmay-replication} & un
solo linaje & comparar ambas fuentes sin cambiar el resto del optimizador & F
\archivo{dwnominate.hpp:55-59, 73}; R 252, 365 & F3, A2 \\
D-01 & SVD con \texttt{DGESDD}, la de wmay; sin LAPACK no se construye & Eigen sustituido en
silencio & la ruta de SVD mueve los óptimos de paneles dinámicos en cientos de nats & F
\archivo{CMakeLists.txt:22-30, 62-68} & F3 \\
D-31 & \texttt{std::sort} por índices y espectral de Eigen, no \texttt{RSORT}, \texttt{RS} ni
\texttt{DSYEV} & port literal & simplicidad, con el umbral de \texttt{REGA}; riesgo: otro orden
de empates & F \archivo{sort_utils.hpp:9-12}, \archivo{simple_ols.hpp:9-12} & F3 \\
D-07 & ciclos fijos sin prueba de convergencia; en el moderno, parada opcional apagada & regla
de parada común & la referencia no tiene criterio de convergencia; se compara a igual
horizonte & F \archivo{dwnominate.cpp:446}; M \archivo{main_cli.cpp:119-122} & F4, T5 \\
D-03 & sin reducción de punto flotante en regiones paralelas; suma en orden de índice &
reducción de OpenMP & otro orden de suma movía los bits bajos e invertía cuasiempates & F
\archivo{likelihood.cpp:423-433} & F4 \\
D-28 & el portal fija \archivo{OMP_NUM_THREADS}; un hilo para Fortran y fiel & hilos del
entorno & antes de D-03, hilos libres dispersaban 29,75 nats; hoy, número declarado y sellado & R 263, 329-336 &
F4, A2 \\
D-11, D-12 & búsquedas NLopt solo en el moderno; motores comparados, no
combinados & NLopt en el fiel; elegir un motor & fidelidad y mejora se responden con motores
distintos & árboles F y M & T5 \\
D-25 & perfil de replicación: tabla de verosimilitud de la referencia, sin cotas globales de
$w_2$ y $\beta$; el portal añade precisión estricta & opciones por omisión & aislamiento
causal: solo cambian las búsquedas & M \archivo{main_cli.cpp:379-390},
\archivo{dwnominate.cpp:569-573}; R 362-364 & T5 \\
D-13 & escalares del moderno en un intervalo local recentrado & caja global & conserva el
alcance de \texttt{SIGMAS} y \texttt{WINT} & M \archivo{parameter_optimizer.hpp:39-42} & T5 \\
D-32 & el moderno aborta si la LL previa es finita y $\Delta\ell<-10^{-8}$ & seguir y exportar &
inferencia: detectar una violación del contrato de bloque & M \archivo{dwnominate.cpp:384-393} & T5 \\
D-06 & grado efectivo: lineal con cinco periodos servidos, cuadrático con seis, cúbico con
siete & grado uniforme & es la regla de la referencia & F \archivo{dwnominate.cpp:1503-1511};
\archivo{wmay.f:1723} & F6 \\
D-10 & exportar en el tiempo normalizado de cada legislador & rejilla global & es donde
optimiza el motor & F \archivo{main_cli.cpp:435-436} & F6 \\
D-26 & sin semilla, el legislador parte del origen & rampa según el padrón & el origen es
neutral; la rampa inyectaba el orden de los identificadores & F
\archivo{csv_loader.cpp:570-609} & F6 \\
D-27 & el portal no siembra en proceso un panel de varios periodos & sembrar con el primer
periodo & dejaría en el origen a quien no sirvió en él, como a 61 de 168 senadores & R 286-299
& F6 \\
D-14 & tres niveles de verificación con evaluador común & coincidencia de salidas & un defecto
de indexación dejó igual la primera dimensión y alejó la verosimilitud &
\archivo{stage_terminal_state.py}, \archivo{map_agreement.py} & F7 \\
D-29 & salida cero, CSV requeridos legibles y directorio inicialmente vacío; excluidos devueltos
& leer lo que haya & rechaza archivos faltantes o previos; no valida toda la cobertura & R 279, 380-391
& A2 \\
D-30 & sello: binario y md5, SVD, hilos, semilla y md5, ciclos, banderas & confiar en la
documentación & hay binarios que no informan su SVD, que pesa más que el mayor desacuerdo entre
motores & R 27-30, 201-207 & F1, A2 \\
\end{longtable}
}

\section{Atributos de calidad: mecanismo y evidencia}
\label{sec:atributos}

El cuadro \ref{tab:atributos} lleva cada atributo de calidad que motiva la reimplementación al
mecanismo que lo persigue y a su evidencia; ninguno se da por demostrado. El rendimiento no se
ha medido todavía, y la carrera P-01 impide afirmar propiedad alguna del motor fiel con más
de un hilo hasta que se corrija.

\begin{table}[tbp]
\centering
\caption{Atributos de calidad: mecanismo de diseño, evidencia y estado.}
\label{tab:atributos}
\footnotesize
\begin{tabular}{L{2.7cm}L{4.2cm}L{2.4cm}L{3.0cm}}
\toprule
Atributo & Mecanismo & Evidencia & Estado \\
\midrule
Mantenibilidad y auditabilidad (RNF09) & configuración explícita; cabeceras por
responsabilidad, dependencias acíclicas; correspondencia uno a uno con la referencia;
desviaciones registradas & F3, A1; D-19, D-22, D-23, D-31 & criterio de diseño; no medido \\
\addlinespace
Operabilidad (RNF03, RNF06) & un portal común a las tres variantes; aceptación con salida cero
y CSV requeridos legibles; excluidos devueltos; sello & F1, A2; D-27
a D-30; anexo \ref{anexo:construccion} & en uso; reejecución desde un clon limpio pendiente
(F1 del backlog) \\
\addlinespace
Rendimiento (RNF04) & paralelismo por votación y legislador; comparación a un hilo; mejora
aceptada solo con resultados idénticos & F4; D-03, D-20, D-21, D-28 & \textbf{no medido}; con
más de un hilo, bloqueado por P-01 \\
\addlinespace
Fidelidad y reproducibilidad (RNF01, RNF02) & \texttt{DGESDD}, aritmética estricta,
\texttt{GRID} apagado, perfil de replicación; ciclos igualados; tres niveles de verificación &
F3, T5, F7; D-01, D-02, D-07, D-14, D-24, D-25 & umbral pendiente del propietario; sin caso de
prueba por hilos \\
\addlinespace
Portabilidad (RNF05) & C++ estándar con CMake; dependencias que detienen la construcción si
faltan; sin código específico de la máquina & D-01, D-02; anexo \ref{anexo:construccion} &
clon limpio en la Etapa 4 \\
\bottomrule
\end{tabular}
\end{table}

\section{Diseño que las figuras no sustituyen}
\label{sec:cuadros}

\subsection{Verificación por niveles}
\label{sec:niveles}

El cuadro \ref{tab:niveles} completa los tres niveles de RF08 con su referente. La
inicialización por mínimos cuadrados de los legisladores es la única rutina portada que el
nivel de componente aún no audita. El conjunto de falsación reintroduce un defecto real ya
corregido y exige que la decisión de aceptación lo rechace, no cada métrica aislada (D-15). Cuatro controles preceden a toda
comparación: paridad de entradas, determinismo por hilos, configuración igualada y tolerancia
declarada entre máquinas.

\begin{table}[tbp]
\centering
\caption{Los tres niveles de verificación de RF08: referente, mecanismo y medida.}
\label{tab:niveles}
\footnotesize
\begin{tabular}{L{2.0cm}L{1.7cm}L{2.8cm}L{3.2cm}L{2.3cm}}
\toprule
Nivel & Referente & Qué se compara & Mecanismo & Medida y aceptación \\
\midrule
Componente & original 2004; las pruebas unitarias, sin referente & cada rutina portada contra
el estado intermedio: dirección de corte, errores por votación, escalares del primer ciclo &
volcado de fases del fiel frente al registro por votación de 2004; en el moderno, bloques
fijados (\archivo{--fix-*}, \archivo{--evaluate-only}) y pruebas unitarias; el fiel declara
esos modos, pero su CLI los deja apagados & tablas por rutina a ciclos igualados; escalares
idénticos tras el primer ciclo \\
\addlinespace
Estado & wmay & la verosimilitud de un estado terminal bajo un único evaluador & estado
preparado; wmay en modo \texttt{evaluate}, cero ciclos & $\ell/N$ y desacuerdo por votación;
umbral $\tau$ por fijar \\
\addlinespace
Mapa & wmay & dos configuraciones emparejadas por (legislador, periodo) & alineación ortogonal
sin reescalar, en marco global y por periodo & $r_1$, $r_2$, $A$, distancia, reflexiones;
cláusulas por fijar (D-15) \\
\bottomrule
\end{tabular}
\end{table}

\subsection{Referentes y protocolos de comparación}
\label{sec:protocolos}

La indicación del 5 de octubre fue comparar con el Fortran original (00:54:23); el profesor
guía no identificó un hash en esa conversación. Los autores la vinculan a la distribución
2004 y al ejecutable reconstruido registrado en el anexo \ref{anexo:procedencia}. Se separan
dos protocolos, porque la referencia del artículo JCC2026 es wmay en sus brazos estáticos y
dinámicos. No se reinterpretan los resultados publicados como resultados de 2004.

\begin{table}[tbp]
\centering
\caption{Protocolos de referencia y alcance de sus medidas.}
\label{tab:protocolos}
\footnotesize
\begin{tabular}{L{2.4cm}L{4.8cm}L{5.2cm}}
\toprule
Propósito & Brazos y configuración & Evaluación y límite \\
\midrule
Reproducción del artículo & wmay y fiel en modo wmay: DGESDD, GRID apagado;
moderno como intervención controlada & puntuación común con wmay (F7), mapas y LL nativa
identificada; no constituye aceptación contra 2004 \\
\addlinespace
Aceptación de tesis E3/E4 & original 2004 (LSVRR, GRID) y fiel con GRID activo;
documentar DGESDD y desviaciones D-19 & mismos paneles, orden, semillas y precisión;
estado, mapas y LL nativa por evaluador. Puntuar ambos con wmay es un control adicional \\
\addlinespace
Rendimiento & original 2004 como referencia externa; fiel antes/después de cada
cambio de memoria & un núcleo, misma máquina, entradas y ciclos; builds y banderas
comparables. La regresión exige salidas idénticas del fiel \\
\bottomrule
\end{tabular}
\end{table}

El modo \archivo{faithful2004-safeguard} conserva GRID, pero no acredita equivalencia de todo
el motor: SVD, exportación y desviaciones deben registrarse. Las comparaciones original/wmay
retenidas del 6 de octubre establecen acuerdo cercano bajo semillas comunes redondeadas;
las recompilaciones reprodujeron cada salida estadounidense propia. No son una medición de
velocidad ni demuestran igualdad entre motores. El umbral RNF01 y su ratificación siguen
pendientes; E2 especifica el protocolo, E3/E4 lo ejecutan.

\subsection{Diseño experimental de las pruebas exhaustivas}
\label{sec:experimental}

Las pruebas exhaustivas de la Etapa 4 establecen, para cada panel de un conjunto fijado de
antemano, si el motor fiel satisface el protocolo original 2004 dentro de la tolerancia de RNF01
\parencite[cap.~5, §2.2.1]{swebok2025}. El diseño, escrito antes de correrlas (tarea H8),
fija paneles, estadísticos de los tres niveles, identidad de las salidas numéricas con uno, dos, cuatro y
doce hilos, y ajustes igualados (anexo \ref{anexo:paneles}). El umbral se propone en su forma:
un $\tau$ sobre la diferencia de verosimilitud por voto, que a diferencia de la absoluta viaja
entre paneles de tamaño distinto, más dos cláusulas de mapa que el conjunto de falsación mostró
sensibles a un defecto real, cada una entre el valor de una versión defectuosa y el de una
corregida (D-15). El valor es decisión del propietario; el piso de ruido de la configuración
igualada, que se medirá en la Etapa 3, lo acota desde abajo.

\subsection{Mejora de rendimiento del motor fiel}
\label{sec:rendimiento}

En la comparación publicada, a un hilo, el motor fiel es más lento que wmay en seis de siete
paneles \parencite{nieves2026jcc}, y el propietario pidió alcanzar el rendimiento a un núcleo
del original 2004 revisando punteros y estructuras de datos. La comparación publicada con
wmay es antecedente; el referente externo de E3/E4 es el original 2004 del cuadro
\ref{tab:protocolos}. Se diseña como prueba de regresión más
prueba de rendimiento \parencite[cap.~5, §2.2.5 y §2.2.7]{swebok2025}; nada está implementado
ni se anticipa magnitud alguna. La hipótesis, que el perfilado debe contrastar, es que el
tiempo se va en recorrer celdas vacías: los votos se guardan con una fila por legislador y
periodo y una columna por votación del panel entero. Los cambios candidatos restringen los
lazos al bloque de cada periodo, guardan los votos contiguos y suprimen trabajo sin efecto en
el resultado (M-01 a M-09). Se acepta un cambio solo si las salidas son idénticas a uno y a
varios hilos (D-20); el protocolo corrige antes P-01, reconstruye la línea base, perfila,
repite corridas intercaladas y aplica los cambios de a uno (D-21; anexo
\ref{anexo:rendimiento}).

\subsection{Diseño del estudio chileno: bootstrap y periodización}
\label{sec:algoritmos}

El estudio chileno, que la sección \ref{sec:propuesta-chile} propone leer como la validación
en condiciones operativas reales, estima el desplazamiento de las medianas de partido en la
primera dimensión durante el 55.º periodo legislativo, con cortes fijados desde el calendario e
inferencia por bootstrap paramétrico \parencite{lewis2004measuring,carroll2009measuring} con
valores iniciales regenerados en cada réplica y ajuste de \textcite{holm1979simple} (RF10,
RF11; algoritmo \ref{alg:bootstrap}). La periodización, que mueve el resultado (RF12), se trata
como parámetro declarado (tarea E3): el diseño principal, pendiente de ratificación (D-16), es
una ventana igualada de doce unidades de cuatro meses en un marco, con tres diseños de
sensibilidad y una regla de discretización de tres cláusulas; la inferencia se hace solo en la
primera dimensión (D-17). El diseño completo, con el cuadro de periodizaciones, está en el
anexo \ref{anexo:chile}.

\subsection{Trazabilidad de los requisitos funcionales al diseño}
\label{sec:trazabilidad}

El cuadro \ref{tab:traza} asigna cada requisito funcional al elemento de diseño que lo
satisface, con lo que queda a la implementación \parencite[cap.~11, §3.4]{swebok2025}, como
revisión del diseño contra los requisitos \parencite[cap.~3, §6.5]{swebok2025}; los no
funcionales están en el cuadro \ref{tab:atributos}.

\begin{table}[tbp]
\centering
\caption{Trazabilidad entre requisitos funcionales y elementos de diseño.}
\label{tab:traza}
\footnotesize
\begin{tabular}{L{1.5cm}L{6.6cm}L{4.3cm}}
\toprule
Req. & Elemento de diseño (dónde se ve) & Pendiente \\
\midrule
RF01, RF02 & validación en la carga y en el portal; dos contratos de semilla (F6, A2; D-26,
D-27) & filas cortas, identificadores no numéricos, códigos fuera de rango (A4) \\
RF03, RF04 & misma maquinaria para uno o varios periodos (F4, F6; D-06) & sin brecha de diseño; pruebas E3 \\
RF05 & módulos uno a uno con la referencia; dos linajes (F3, A1; D-22 a D-24) & ninguno \\
RF06 & variante en las cuatro búsquedas; perfil de replicación (T5; D-25) & registrar la SVD en
el resumen del moderno \\
RF07 & marco de exportación declarado (F6; D-10) & A5; resumen y exportación 1D (I2) \\
RF08, RF09 & tres niveles; desacuerdo por votación (F7; cuadro \ref{tab:niveles}; D-14) &
OLS; modos de componente (I4); claves y productor RF09 (I5) \\
RF10, RF11 & bootstrap con semillas regeneradas, Holm, elegibilidad (alg. \ref{alg:bootstrap})
& brazo fiel sobre el diseño principal \\
RF12 & periodización declarada (cuadro \ref{tab:disenos}; D-16) & rejilla de igual recuento \\
\bottomrule
\end{tabular}
\end{table}
